
\input{$layout}

\enableregime[utf]
\mainlanguage[cn]
\language[cn]
\setscript[hanzi] % hyphenation

%%%%%%%
% 字体
\usetypescriptfile[type-imp-myfonts]

% mode 由命令行调用时提供
\startmode [tc]
  \usetypescript[tcfonts]
  \setupbodyfont[tcfonts,cg]
\stopmode

\startmode [sc]
  \usetypescript[scfonts]
  \setupbodyfont[scfonts,cg]
\stopmode
% 字体
%%%%%%%

\startmode [debug]
  \showframe
\stopmode




\setupindenting[yes,first,2em]
\setupheads[indentnext=yes]

\usecolors[crayola]
%\usecolors[ral]


%%%%%%%%%%%%
% 页脚页眉
%\setupheader
%  [text]
%  [before={\startframed[frame=off,bottomframe=on,framecolor=black,]},
%   after={\stopframed},
%  ]
%\setupfooter
%  [text]
%  [before={\startframed[frame=off,bottomframe=off,topframe=on,framecolor=black,]},
%  after={\stopframed},
%  ]

%%%%%%%%%%%%%%%%%%%%%%%%%%%%%%%%%%%%%%%%%%%%%%%%%%%%%%%%%%%%%%%%%%%%%%%%%%%%%%%%%%%%%%%%%%%%%%%%%%%%%%%%%%%%%%%%%%%%%%%%
%                                          Header and Footer
%
%\setupheader
%  [text]
%  [before={\startframed[frame=off,bottomframe=on,topframe=off,framecolor=black,]},
%   after={\stopframed},
%  ]
%\setupfooter
%  [text]
%  [before={\startframed[frame=off,bottomframe=off,topframe=on,framecolor=black,]},
%   after={\stopframed},
%  ]


\startsetups mainheaderfooter
\setupheadertexts[]
%\setupheadertexts[{\tf\bf\ss \somenamedheadnumber{chapter}{current}. \getmarking[chapter]}][]
\setupfootertexts[{\it \pagenumber}]
\stopsetups


\startsetups clearheaderfooter
\setupheadertexts[]
\setupfootertexts[]
\stopsetups
%
%
%%%%%%%%%%%%%%%%%%%%%%%%%%%%%%%%%%%%%%%%%%%%%%%%%%%%%%%%%%%%%%%%%%%%%%%%%%%%%%%%%%%%%%%%%%%%%%%%%%%%%%%%%%%%%%%%%%%%%%%%

% 页脚页眉
%%%%%%%%%%%


%%%%%%%%%%%%%%%
% PDF 属性信息
\setupinteraction[focus=standard] % 重要功能，添加上后才能链接到位置
\setupinteraction[state=start,style=normal]
\setupinteraction[
    title={$title},
    author={$author},
    keyword={$keyword},
    creator={https://github.com/meng89/nikaya},
]


% rulecolor：正文与脚注之间的隔离线条颜色。textcolor：正文注解引号的颜色
\setupfootnotes[textcolor=blue]
% headcolor：脚注引号的颜色。color：脚注文本的颜色
\setupnotation[footnote][way=bypage,numbercommand={\tfx},distance=0.1em,headcolor=blue] % note text is blue

\def\zhfootnote#1{\footnote{\setscript[hanzi]\rm \tfx #1}}


% PDF 属性信息
%%%%%%%%%%%%%%%


%%%%%%%%%%%%
% PDF书签
%\definelist[tobias]

%\placebookmarks[appendix,part,chapter,section,subsection,tobias][chapter][number=no]
\definehead[notitle][part]
\definehead[nolist][part]

\definelist[notitle][part]

\placebookmarks[part,title,subject,subsubject,subsubsubject,subsubsubsubject,subsubsubsubsubject,subsubsubsubsubsubject,subsubsubsubsubsubsubject,notitle,nolist]

\setupcombinedlist[content]
    [list={notitle,part,title,subject,subsubject,subsubsubject,subsubsubsubject,subsubsubsubsubject,subsubsubsubsubsubject,subsubsubsubsubsubsubject},
     alternative=c,
    ]

\setuplist[part,notitle][margin=0em]
\setuplist[title][margin=1em]
\setuplist[subject][margin=2em]
\setuplist[subsubject,][margin=3em]
\setuplist[subsubsubject,][margin=4em]
\setuplist[subsubsubsubject,][margin=5em]
\setuplist[subsubsubsubsubject,][margin=6em]
\setuplist[subsubsubsubsubsubject,][margin=7em]
\setuplist[subsubsubsubsubsubsubject][margin=8em]


\setuphead[part][style=\tfd\bf\ss]
\setuphead[title][style=\tfc\bf\ss]
\setuphead[subject][style=\tfb\bf\ss]
\setuphead[subsubject][style=\tfa\bf\ss]
\setuphead[subsubsubject][style=\tf\bf\ss]
\setuphead[subsubsubsubject][ style=\tfx\bf\ss]
\setuphead[subsubsubsubsubject][style=\tfxx\bf\ss]
\setuphead[subsubsubsubsubsubject][style=\tfxx\bf\ss]
\setuphead[subsubsubsubsubsubsubject][ style=\tfxx\bf\ss]

\setuphead[
  notitle,
  part,
  title,
  subject,
  subsubject,
  subsubsubject,
  subsubsubsubject,
  subsubsubsubsubject,
  subsubsubsubsubsubject,
  subsubsubsubsubsubsubject]
  [
    number=no,
    placehead=yes,
    page=no,
    incrementnumber=list,
    textdistance=0em,
    color=darkblue,
    topframe=off,
    before=, %before={\startalignment[middle]before},%before={\blank[halfline]},
    after=, %\stopalignment},%after={\blank[halfline]},
    %align=middle,
    aftersection={\blank[3*halfline]},
    interaction=list,
    alternative=text,
]

\setuphead[notitle][
  placehead=no,
  interaction=list,
  page=yes,
]
\setuphead[nolist][
  interaction=list,
  page=yes,
  align=middle,
  alternative=normal,
]




\setupinteractionscreen[option=bookmark]
\enabledirectives[references.bookmarks.preroll]

\enabledirectives[backend.pdf.fixhighlight]


\definecolor[pdfhighlight:莊春江][r=.93,g=.93,b=.93,]



\useexternalfigure[cover][$cover_image]
\defineoverlay[cover][{\externalfigure[cover][width=\paperwidth]}]

\startsetups Background:cover
\setupbackgrounds[page][background=cover]
\stopsetups

\startsetups Background:off
% 只要把 background 留空或这设置成无效的名字，就会没有背景
\setupbackgrounds[page][background=,]
\stopsetups

%\showframe

%%%%%%%%%%%%%%%%%%%%%%%%%%%%%%%%%%%%%%%%%%%%%%%%%%%%%%%%%%%%%%%%%%%%%%%%%%%%%%%%%%%%%%%%%%%%%%%%%%%%%%%%%%%%%%%%%%%%%%%%

\startdocument

\setups[clearheaderfooter]
%\setuppagenumber[number=1]
\setups[Background:cover]
%\bookmark[notitle]{封面}
\null
\page
\setups[Background:off]


%
\startnolist[title={目录},bookmark={页面目录},list=,]
\startcolor[TropicalRainForest]
\placecontent
\stopcolor
\page[yes]
\stopnolist

\bookmark[mylist]{$line1}
\vfill % 在内容前添加弹性空间
\startalignment[center]
\tfd{$line1}

\tfd{$line2}
\stopalignment
\vfill % 在内容后添加弹性空间
\page

\setups[mainheaderfooter]
\input{suttas.tex}
\page[yes]

\setups[clearheaderfooter]
\setupindenting[no]
\setupheads[indentnext=no]
\input{readme.tex}

\stopdocument
